\section{Module \texttt{window.c}}

Le module \texttt{window.c} construit les fenêtres principales de l'application. Il encapsule la création d'un \texttt{GtkApplicationWindow}, l'application des propriétés usuelles et la gestion d'un fond personnalisé par CSS.

\subsection{Fonction \texttt{apply\_window\_background}}

Cette fonction privée applique un style CSS ciblé sur une seule fenêtre. L'idée est de nommer le widget puis de générer dynamiquement une règle CSS qui peut définir soit une couleur de fond seule, soit une couleur combinée à une image.

\begin{lstlisting}[language=C]
static void apply_window_background(const window_config *config, GtkWidget *window) {
    if (config == NULL || window == NULL || config->background_color == NULL || config->background_color[0] == '\0') {
        return;
    }

    const char *name = config->widget_name;
    if (name == NULL || name[0] == '\0') {
        name = "wrapper-window";
        gtk_widget_set_name(window, name);
    }

    GtkCssProvider *provider = gtk_css_provider_new();
    char *css;

    if (config->background_image_path != NULL && config->background_image_path[0] != '\0') {
        css = g_strdup_printf(
            "window#%s, window#%s > * {"
            "background-color: %s;"
            "background-image: url('%s');"
            "background-size: cover;"
            "background-position: center;"
            "}",
            name,
            name,
            config->background_color,
            config->background_image_path
        );
    } else {
        css = g_strdup_printf(
            "window#%s, window#%s > * { background-color: %s; }",
            name,
            name,
            config->background_color
        );
    }

    gtk_css_provider_load_from_string(provider, css);
    gtk_style_context_add_provider_for_display(
        gtk_widget_get_display(window),
        GTK_STYLE_PROVIDER(provider),
        GTK_STYLE_PROVIDER_PRIORITY_APPLICATION
    );

    g_free(css);
    g_object_unref(provider);
}
\end{lstlisting}

\subsection{Fonction \texttt{create\_window}}

La fonction principale du module instancie la fenêtre, applique les propriétés déclarées dans \texttt{window\_config}, délègue la personnalisation visuelle à \texttt{apply\_window\_background} puis applique le style commun partagé avec les autres widgets.

\begin{lstlisting}[language=C]
GtkWidget *create_window(GtkApplication *app, const window_config *config) {
    GtkWidget *window = gtk_application_window_new(app);
    if (config == NULL) {
        return window;
    }

    if (config->title) {
        gtk_window_set_title(GTK_WINDOW(window), config->title);
    }
    if (config->icon_name) {
        gtk_window_set_icon_name(GTK_WINDOW(window), config->icon_name);
    }
    if (config->widget_name) {
        gtk_widget_set_name(window, config->widget_name);
    }

    gtk_window_set_default_size(GTK_WINDOW(window), config->default_width, config->default_height);
    gtk_window_set_resizable(GTK_WINDOW(window), config->resizable);
    gtk_window_set_decorated(GTK_WINDOW(window), config->decorated);
    gtk_window_set_modal(GTK_WINDOW(window), config->modal);

    if (config->min_width > 0 || config->min_height > 0) {
        gtk_widget_set_size_request(window, config->min_width, config->min_height);
    }

    apply_window_background(config, window);
    apply_widget_style(window, &config->style);

    if (config->maximized) {
        gtk_window_maximize(GTK_WINDOW(window));
    }
    if (config->present_on_create) {
        gtk_window_present(GTK_WINDOW(window));
    }

    return window;
}
\end{lstlisting}

\subsection{APIs GTK utilisées}

\begin{itemize}
    \item \textbf{\texttt{gtk\_application\_window\_new()}} : Crée une fenêtre principale intégrée au cycle de vie d'une \texttt{GtkApplication}. Le module l'utilise pour bénéficier de la gestion automatique des actions et de la fermeture de l'application.
    \item \textbf{\texttt{gtk\_widget\_set\_name()}} : Assigne un identifiant unique (ID) au widget, utilisable comme sélecteur dans une feuille de style CSS. Le wrapper s'en sert pour cibler spécifiquement la fenêtre lors de l'application dynamique de thèmes personnalisés.
    \item \textbf{\texttt{gtk\_css\_provider\_new()}} : Instancie un objet capable de lire et d'analyser des données CSS. Elle permet au module de transformer les paramètres de couleur et d'image de fond du wrapper en règles de style compréhensibles par le moteur de rendu de GTK.
    \item \textbf{\texttt{gtk\_style\_context\_add\_provider\_for\_display()}} : Attache un fournisseur de style à l'ensemble de l'écran (display). Cette API est cruciale pour que les styles générés dynamiquement par le wrapper soient pris en compte par le moteur CSS de GTK lors du rendu de la fenêtre.
    \item \textbf{\texttt{gtk\_window\_set\_default\_size()}} : Détermine les dimensions d'ouverture par défaut de la fenêtre. Elle est utilisée pour initialiser l'espace de travail de l'utilisateur selon les préconisations fournies dans la configuration.
    \item \textbf{\texttt{gtk\_window\_present()}} : Rend la fenêtre visible et demande au gestionnaire de fenêtres de lui donner le focus. Le wrapper l'appelle systématiquement en fin de création pour garantir que l'utilisateur voit immédiatement l'interface générée.
\end{itemize}

\newpage
